\documentclass[11pt]{article}
\usepackage[T1]{fontenc}
\usepackage{lmodern}
\usepackage[margin=1in]{geometry}
\usepackage{amsmath,amssymb}
\usepackage[hidelinks]{hyperref}
\usepackage{xcolor}
\setlength{\parindent}{0pt}
\setlength{\parskip}{0.7em}

\title{Literature Status: Bourgain's Slicing Problem}
\author{Run \texttt{fa\_banach\_001} -- agent \texttt{agent\_lane\_02}}
\date{19 July 2026}

\begin{document}
\maketitle

\begin{center}
\fcolorbox{black}{gray!10}{\parbox{0.88\linewidth}{
\textbf{Status: literature\_already\_answered.}
The source question was affirmatively resolved in separate later literature.
This packet records an exact literature match, not a new proof.
}}
\end{center}

\section*{Source question}

Paouris and Werner, \emph{Relative entropy of cone measures and
$L_p$ centroid bodies}, arXiv:0909.4361v1 (2009), Section 2, PDF page 6,
state:
\begin{quote}
``A major open problem in convex geometry asks if there exists a universal
constant $C>0$ such that $L_K\leq C$.''
\end{quote}
Here $L_K$ is the affine-invariant isotropic constant of a convex body $K$.
The source records the then-best estimate $L_K\leq Cn^{1/4}$.

\section*{Supporting answer}

Klartag and Lehec, \emph{Affirmative Resolution of Bourgain's Slicing
Problem using Guan's Bound}, arXiv:2412.15044v1 (2024), published in
\emph{Geometric and Functional Analysis} 35 (2025), 1147--1168, explicitly
identify and resolve the same problem. On PDF page 2 they define
\[
  L_n=\sup\{L_K:K\subseteq\mathbb{R}^n\}
\]
and prove in Theorem 1.2 that
\[
  \sup_{n\geq 1} L_n<\infty.
\]
Equivalently, Theorem 1.1 proves that for every volume-one convex body
$K\subseteq\mathbb{R}^n$, some hyperplane $H$ satisfies
$\operatorname{Vol}_{n-1}(K\cap H)>c$ for a universal $c>0$.

\newpage
\section*{Identification}

The quantifiers match exactly: boundedness of $\sup_n L_n$ says that one
constant bounds $L_K$ for every dimension and every convex body. The
supporting theorem is stronger than needed for the symmetric context around
the source passage. The supporting authors explicitly know they are solving
Bourgain's slicing problem, as stated in their title, abstract, introduction,
and Theorems 1.1--1.2.

\section*{Scope limitations}

The universal-boundedness question in the 2009 source is fully answered.
Stronger extremizer conjectures mentioned in the supporting paper -- simplex
maximality in the general case and cube maximality in the centrally symmetric
case -- are separate questions and are not claimed resolved here.

\section*{Search evidence and recommendation}

The four lightweight run indexes contained no prior record for arXiv:0909.4361
or this exact source-to-answer identification. Related packets for
arXiv:math/0512207 and arXiv:1310.1204 already record that arXiv:2412.15044
answers the same named conjecture in other source papers; this packet adds
only source-specific status memory. The supporting PDF was checked at pages
1--2 and Theorems 1.1--1.2. The authors' publication record confirms its 2025
GAFA publication.

Human review should check only the normalization match between source PDF
page 6 and supporting PDF page 2. This lightweight retrieval packet does not
purport to audit the published proof.

\begin{thebibliography}{9}
\bibitem{PW}
G. Paouris and E. M. Werner,
\emph{Relative entropy of cone measures and $L_p$ centroid bodies},
arXiv:0909.4361v1 (2009),
\url{https://arxiv.org/abs/0909.4361}.

\bibitem{KL}
B. Klartag and J. Lehec,
\emph{Affirmative Resolution of Bourgain's Slicing Problem using Guan's Bound},
Geom. Funct. Anal. \textbf{35} (2025), 1147--1168;
arXiv:2412.15044v1,
\url{https://doi.org/10.1007/s00039-025-00718-w}.
\end{thebibliography}

\end{document}
